\begin{center}
    {\LARGE \textbf{2024无机化学与物理化学}} \\[2ex]
    {\large 时量180分钟 \quad 满分150分}
\end{center}

\vspace{0.5cm}
\noindent\textbf{注意事项:}
\begin{enumerate}[label=\arabic*., parsep=0pt, itemsep=0pt]
    \item 答题前,考生先将自己的姓名、学号、考场号、座位号填写在答题卡上,将条形码准确粘贴在考生信息条形码粘贴区。
    \item 考试结束后,将本试卷与答题纸一并收回。
    \item 回答各个问题时,将答案写在答题纸上,写在本试卷上无效。
\end{enumerate}

\vspace{1cm}
\begin{center}
    {\Large \textbf{无机化学部分(共75分)}}
\end{center}

\vspace{0.5cm}
\noindent\textbf{一、选择题(20 pts.)}

\begin{enumerate}[label=\arabic*., itemsep=1.5ex]
    \item 下列物质能稳定存在的是
    \begin{tasks}(4)
        \task $\ce{SnCl4}$
        \task $\ce{GeCl4}$
        \task $\ce{PbCl4}$
        \task $\ce{SbCl5}$
    \end{tasks}

    \item 由 $\ce{Cr2O3}$ 制备铬酸盐所需要用到的氧化剂为
    \begin{tasks}(4)
        \task $\ce{KClO3 + KCl}$
        \task $\ce{MnO2}$
        \task $\ce{H2O2}$
        \task $\ce{HNO3 + HCl}$
    \end{tasks}

    \item 下列物质中，酸性最强的是
    \begin{tasks}(4)
        \task $\ce{H2S}$
        \task $\ce{H2SO3}$
        \task $\ce{H2SO4}$
        \task $\ce{H2S2O7}$
    \end{tasks}

    \item 制备 $\ce{F2}$ 实际所采用的方法是
    \begin{tasks}(4)
        \task 电解 $\ce{HF}$
        \task 电解 $\ce{CaF2}$
        \task 电解 $\ce{KHF2}$
        \task 电解 $\ce{NH4F}$
    \end{tasks}

    \item 与其他化合物颜色明显不同的是
    \begin{tasks}(4)
        \task $\ce{Cu(OH)2*CuCO3}$
        \task $\ce{NaCrO2}$
        \task $\ce{K2MnO4}$
        \task $\ce{K2CrO4}$
    \end{tasks}

    \item 既易溶于稀氢氧化钠溶液，又易溶于氨水的是
    \begin{tasks}(4)
        \task $\ce{Cu(OH)2}$
        \task $\ce{Ag2O}$
        \task $\ce{Zn(OH)2}$
        \task $\ce{Cd(OH)2}$
    \end{tasks}

    \item 下列分子中，不存在 d-pπ 配键的是
    \begin{tasks}(4)
        \task $\ce{SOCl2}$
        \task $\ce{HNO3}$
        \task $\ce{H3PO4}$
        \task $\ce{HClO2}$
    \end{tasks}

    \item 下列分子或离子的键能排序正确的是
    \begin{tasks}(2)
        \task $\ce{N2^{+}} > \ce{N2} > \ce{N2^{-}}$
        \task $\ce{O2^{+}} > \ce{O2} > \ce{O2^{-}}$
        \task $\ce{H2} > \ce{F2} > \ce{N2}$
        \task $\ce{N2} > \ce{O2} > \ce{O2^{+}}$
    \end{tasks}

    \item 将 $\ce{SiF4}$ 通入 $\ce{NaOH}$ 溶液中，主要产物是
    \begin{tasks}(4)
        \task $\ce{H4SiO4 + NaF}$
        \task $\ce{Na2SiO3 + NaF}$
        \task $\ce{Na2SiO3 + Na2SiF6}$
        \task $\ce{SiO2 + HF}$
    \end{tasks}

    \item 下列配位化合物中，可以由简单盐溶液与 $\ce{KCN}$ 溶液直接得到的是
    \begin{tasks}(4)
        \task $\ce{K3[Fe(CN)6]}$
        \task $\ce{K4[Co(CN)6]}$
        \task $\ce{K2[Ni(CN)4]}$
        \task $\ce{K2[Cu(CN)4]}$
    \end{tasks}
\end{enumerate}

\vspace{0.5cm}
\noindent\textbf{二、简答题(45 pts.)}

\begin{enumerate}[label=\arabic*., itemsep=1.5ex]
    \item 解释下列实验现象
    \begin{enumerate}[label=(\arabic*), noitemsep, topsep=0pt, partopsep=0pt]
        \item 在煤气灯上加热含结晶水的硝酸铜晶体，观察到黑色固体和红棕色气体。
        \item 向氯化钴溶液中滴加氨水，先产生蓝绿色沉淀，氨水过量时沉淀又溶解。
        \item 向含有碘酸钾的淀粉水溶液中滴加亚硫酸钠，溶液变为蓝色，亚硫酸钠过量时蓝色褪去。
    \end{enumerate}

    \item 用价键理论解释为什么 $[\ce{Ni(NH3)6}]^{2+}$ 为六配位的八面体构型，$[\ce{Ni(CN)4}]^{2-}$ 为四配位的正方形构型。

    \item 分离 $\ce{Al^{3+}}$、$\ce{Zn^{2+}}$、$\ce{Cr^{3+}}$、$\ce{Cu^{2+}}$。

    \item 用化学法由二氧化锰制备高锰酸钾。

    \item 有三瓶白色固体试剂，分别为 $\ce{NaCl}$，$\ce{NaBr}$，$\ce{NaI}$，试用三种方法鉴别。

    \item 解释为什么酸性条件下高锰酸钾和亚硝酸钠反应溶液褪色，碱性条件下高锰酸钾和亚硝酸钠反应溶液变绿。

    \item 在 $\ce{KCl}$ 中混有 $\ce{K2SO4}$，除杂。
\end{enumerate}

\vspace{0.5cm}
\noindent\textbf{三、推断题(10 pts.)}

无色晶体A易溶于水。A的溶液与酸性 $\ce{KI}$ 溶液作用溶液变黄，说明有B生成。A的溶液煮沸一段时间后加入 $\ce{BaCl2}$ 溶液有白色沉淀C生成，C不溶于强酸。A的溶液与 $\ce{MnSO4}$ 混合后加稀硝酸和几滴 $\ce{AgNO3}$ 溶液，加热，溶液变红，说明有D生成。A溶液与 $\ce{KOH}$ 混合后加热，有气体E放出。气体E通入 $\ce{AgNO3}$ 溶液则有棕黑色沉淀F生成，E通入过量则沉淀溶解为无色溶液。

试给出A、B、C、D、E和F所代表的物质的化学式，并用化学反应方程式表示 A→B，A→D，A→E，以及生成 F 的过程。

\vspace{1cm}
\begin{center}
    {\Large \textbf{物理化学部分(共75分)}}
\end{center}

\vspace{0.5cm}
\noindent\textbf{四、选择题(24 pts.)}

\begin{enumerate}[label=\arabic*., start=1, itemsep=1.5ex]
    \item 下述说法哪一种不正确？
    \begin{tasks}(1) 
        \task 理想气体经绝热自由膨胀后，其内能变化为零
        \task 非理想气体经绝热自由膨胀后，其内能变化不一定为零
        \task 非理想气体经绝热膨胀后，其温度一定降低
        \task 非理想气体经一不可逆循环，其内能变化为零
    \end{tasks}

    \item 下述方法中，哪一种对于消灭蚂蟥比较有效？
    \begin{tasks}(4)
        \task 击打
        \task 刀割
        \task 晾晒
        \task 撒盐
    \end{tasks}

    \item 在碰撞理论中校正因子 $P$ 小于 1 的主要因素是：
    \begin{tasks}(2)
        \task 反应体系是非理想的
        \task 空间的位阻效应
        \task 分子碰撞的激烈程度不够
        \task 分子间的作用力
    \end{tasks}

    \item 298 K 时两个级数相同的反应 I、II，活化能 $E_\text{I} = E_\text{II}$，若速率常数 $k_\text{I} = 10 k_\text{II}$，则两反应之活化熵相差：
    \begin{tasks}(2)
        \task $0.6\ \text{J}\cdot\text{K}^{-1}\cdot\text{mol}^{-1}$
        \task $10\ \text{J}\cdot\text{K}^{-1}\cdot\text{mol}^{-1}$
        \task $19\ \text{J}\cdot\text{K}^{-1}\cdot\text{mol}^{-1}$
        \task $190\ \text{J}\cdot\text{K}^{-1}\cdot\text{mol}^{-1}$
    \end{tasks}

    \item 一恒压反应体系，若产物与反应物的 $\Delta C_p > 0$，则此反应
    \begin{tasks}(4)
        \task 吸热
        \task 放热
        \task 无热效应
        \task 吸放热不能肯定
    \end{tasks}

    \item 理想气体反应 $\ce{CO(g) + 2H2(g) = CH3OH(g)}$ 的 $\Delta_\text{r} G_\text{m}^\ominus$ 与温度 $T$ 的关系为：\\
    $\Delta_\text{r} G_\text{m}^\ominus / \text{J}\cdot\text{mol}^{-1} = -21660 + 52.92(T/\text{K})$，若使在标准状态下的反应向右进行，则应控制反应的温度
    \begin{tasks}(2)
        \task 必须高于 409.3 K
        \task 必须低于 409.3 K
        \task 必须等于 409.3 K
        \task 必须低于 $409.3\ ^\circ\text{C}$
    \end{tasks}

    \item 金属活泼性排在 $\ce{H2}$ 之前的金属离子，如 $\ce{Na+}$ 能优先于 $\ce{H+}$ 在汞阴极上析出，这是由于
    \begin{tasks}(1)
        \task $\phi^\ominus(\ce{Na+ / Na}) < \phi^\ominus(\ce{H+ / H2})$
        \task $\eta(\ce{Na}) < \eta(\ce{H2})$
        \task $\phi(\ce{Na+ / Na}) < \phi(\ce{H+ / H2})$
        \task $\ce{H2}$ 在汞上析出有很大的超电势，以至于 $\phi(\ce{Na+ / Na}) > \phi(\ce{H+ / H2})$
    \end{tasks}

    \item 使用瑞利（Reyleigh）散射光强度公式，在下列问题中可以解决的问题是
    \begin{tasks}(2)
        \task 溶胶粒子的大小
        \task 溶胶粒子的形状
        \task 测量散射光的波长
        \task 测量散射光的振幅
    \end{tasks}

    \item 一定量的理想气体从同一始态出发，分别经 (1)等温压缩，(2)绝热压缩到具有相同压力的终态，以 $H_1$、$H_2$ 分别表示两个终态的焓值，则有
    \begin{tasks}(4)
        \task $H_1 > H_2$
        \task $H_1 = H_2$
        \task $H_1 < H_2$
        \task $H_1 \le H_2$
    \end{tasks}

    \item 在还原性酸性溶液中，$\ce{Zn}$ 的腐蚀速度较 $\ce{Fe}$ 为小，其原因是
    \begin{tasks}(2)
        \task $\phi(\ce{Zn^{2+} / Zn})(\text{平}) < \phi(\ce{Fe^{2+} / Fe})(\text{平})$
        \task $\phi(\ce{Zn^{2+} / Zn}) < \phi(\ce{Fe^{2+} / Fe})$
        \task $\phi(\ce{H+ / H2})(\text{平, } \ce{Zn}) < \phi(\ce{H+ / H2})(\text{平, } \ce{Fe})$
        \task $\phi(\ce{H+ / H2})(\ce{Zn}) < \phi(\ce{H+ / H2})(\ce{Fe})$
    \end{tasks}

    \item 已知有下列一组公式可用于理想气体绝热过程功的计算：

    \vspace{1ex}
    \begin{tabular}{p{0.45\textwidth} p{0.45\textwidth}}
        (1) $W = C_V(T_1 - T_2)$ & (2) $[1/(\gamma - 1)](p_1 V_1 - p_2 V_2)$ \\[1ex]
        (3) $[p_1 V_1 / (\gamma - 1)][1 - (V_1/V_2)^{\gamma - 1}]$ & (4) $[p_1 V_1 / (\gamma - 1)][1 - (p_2/p_1)^{1 - 1/\gamma}]$ \\[1ex]
        (5) $[p_1 V_1 / (\gamma - 1)][1 - (p_2 V_2 / p_1 V_1)]$ & (6) $[R/(\gamma - 1)](T_1 - T_2)$
    \end{tabular}
    \vspace{1ex}
    
    但这些公式只适于绝热可逆过程的是
    \begin{tasks}(4)
        \task (1), (2)
        \task (3), (4)
        \task (5), (6)
        \task (4), (5)
    \end{tasks}

    \item 以石墨为阳极，电解 $0.01\ \text{mol}\cdot\text{kg}^{-1}\ \ce{NaCl}$ 溶液，在阳极上首先析出
    \begin{tasks}(2)
        \task $\ce{Cl2}$
        \task $\ce{O2}$
        \task $\ce{Cl2}$ 与 $\ce{O2}$ 混合气体
        \task 无气体析出
    \end{tasks}
    已知：$\phi^\ominus(\ce{Cl2/Cl-}) = 1.36\ \text{V}$，$\eta(\ce{Cl2}) = 0\ \text{V}$，\\
    \hspace*{3em}$\phi^\ominus(\ce{O2/OH-}) = 0.401\ \text{V}$，$\eta(\ce{O2}) = 0.8\ \text{V}$。
\end{enumerate}

\vspace{0.5cm}
\noindent\textbf{五、计算题(10 pts.)}

$1.22\times 10^{-2}\,\text{kg}$ 苯甲酸溶于 $0.10\,\text{kg}$ 乙醇后，使乙醇的沸点升高了 $1.13\,\text{K}$，若将 $1.22\times 10^{-2}\,\text{kg}$ 苯甲酸溶于 $0.10\,\text{kg}$ 苯中，则苯沸点升高 $1.36\,\text{K}$。计算苯甲酸在两种溶剂中的摩尔质量。计算结果说明什么问题。

已知：$k_b(\text{乙醇})=1.19\,\text{K·kg·mol}^{-3}$，$k_b(\text{苯})=2.60\,\text{K·kg·mol}^{-3}$。

\vspace{0.5cm}
\noindent\textbf{六、计算题(15 pts.)}

（1）利用下面数据绘制 $\text{Mg-Cu}$ 体系的相图。$\text{Mg}$（熔点为 $648^\circ\text{C}$）和 $\text{Cu}$（熔点为 $1085^\circ\text{C}$）形成两种化合物：$\text{MgCu}_2$（熔点为 $800^\circ\text{C}$），$\text{Mg}_2\text{Cu}$（熔点为 $580^\circ\text{C}$）；形成三个低共熔混合物，其组成分别为 $w(\text{Mg})=0.10$，$w(\text{Mg})=0.33$，$w(\text{Mg})=0.65$，它们的熔点分别为 $690^\circ\text{C}$，$560^\circ\text{C}$ 和 $380^\circ\text{C}$。已知 $M_r(\text{Cu})=63.55\,\text{g·mol}^{-1}$，$M_r(\text{Mg})=24.31\,\text{g·mol}^{-1}$。

（2）画出组成为 $w(\text{Mg})=0.25$ 的 $900^\circ\text{C}$ 的熔体降温到 $100^\circ\text{C}$ 时的步冷曲线。

（3）当（2）中的熔体 $1\,\text{kg}$ 降温至无限接近 $560^\circ\text{C}$ 时，可得到何种化合物？其质量为若干？

\vspace{0.5cm}
\noindent\textbf{七、计算题(13 pts.)}

电池 $\text{Ag(s)}|\text{AgBr(s)}|\text{HBr(0.1 mol·kg}^{-1})|\text{H}_2(0.01p^\ominus),\text{pt}$，$298\,\text{K}$ 时，$E=0.165\,\text{V}$，当电子得失为 $1\,\text{mol}$ 时，$\Delta_r H_m = -50.0\,\text{kJ·mol}^{-1}$，电池反应平衡常数 $K^\theta = 0.0301$，$E^\ominus(\text{Ag}^+|\text{Ag})=0.800\,\text{V}$，设活度系数均为 $1$。

（1）写出电极与电池反应；
（2）计算 $298\,\text{K}$ 时 $\text{AgBr(s)}$ 的 $K_{\text{sp}}$；
（3）求电池反应的可逆反应热 $Q_R$；
（4）计算电池的温度系数。

\vspace{0.5cm}
\noindent\textbf{八、计算题(13 pts.)}

$\text{NO}$ 高温均相分解是二级反应，反应为：$2\text{NO(g)}\rightarrow\text{N}_2(\text{g})+\text{O}_2(\text{g})$，实验测得 $1423\,\text{K}$ 时速率常数为 $1.843\times 10^{-3}\,\text{dm}^3\cdot\text{mol}^{-1}\cdot\text{s}^{-1}$，$1681\,\text{K}$ 时速率常数为 $5.743\times 10^{-2}\,\text{dm}^3\cdot\text{mol}^{-1}\cdot\text{s}^{-1}$。求：

(1) 等容反应的活化能；
(2) 用过渡态理论计算 $1423\,\text{K}$ 时反应的 $\Delta^\ddagger H_\text{m}^\ominus$、$\Delta^\ddagger U_\text{m}^\ominus$、$\Delta^\ddagger S_\text{m}^\ominus\ (c^\ominus)$ 及 $\Delta^\ddagger G_\text{m}^\ominus\ (c^\ominus)$。